% SPDX-License-Identifier: Apache-2.0
\section{A Choice on a Constant}
\label{sec:x-const-mux}

A unit generic over its configuration may choose between two values on
a constant of the configuration. A \code{mux} or an \code{if} that
yields a value, on the constant itself, folds when the unit is lowered
if the operand it chooses brings its own width: the netlist holds that
operand and no mux, and the input not chosen stays a port of the unit,
read by nothing. A choice on a choice on the constant folds too, since
the inner one folds to a bit.

The other forms stay a mux whose condition is a constant: a selector
worked out of the constant with \code{!}, \code{\&} or \code{==}, and an operand
with no width of its own, a bare number or a sum such as \code{(b + b)},
which takes its width from the other operand. A \code{bool} constant
is a one-bit literal wherever it is written, because its Rust type is
asked when the unit is lowered, so its negation is one bit wide in
Verilog as well and chooses as the run does: \code{q} reads it bare,
under \code{\&\&} and \code{==}, behind a \code{let} and as a
branch of a choice. An \code{if} statement
on the constant is folded, and a check holds the constant. The
netlists compute what the run computes, and nvc accepts every line.
\code{Both} holds the unit twice, once for each value of the flag, and
its netlist is replayed against the trace under nvc and Verilator, so
both lowerings are compiled by both simulators.

\lstinputlisting[caption={\code{ex\_const\_mux.rs}. Run with \code{bazel run //lib/examples:ex\_const\_mux}.},label={lst:x-const-mux}]{lib/examples/ex_const_mux.rs}

\lstinputlisting[style=ascii,caption={What \code{ex\_const\_mux} printed when the build ran it: the outputs of the two configurations, then the lines of each lowering that drive them.},label={lst:x-const-mux-out}]{out_const_mux.txt}
